% figures/w05-current-triangle.svg — 대지속도는 대수속도와 조류의 합이다.
%
%   bash _tools/tikz2svg.sh figures/src/w05-current-triangle.tex figures/w05-current-triangle.svg
%
% 2026-10-17, 사용자의 한국어 참고 그림("해류 삼각형")을 영어로 옮겨 그렸다.
% W05 §5-5 가 쓰던 w04-course-crab 은 조류가 없는 그림이라 이 절의 주제를
% 다 담지 못한다. 이 그림은 각을 **셋이 아니라 다섯**으로 보여 준다.
%
% 각도 규약: 북이 위, 동이 오른쪽. 방위각 a 인 단위벡터의 화면 성분은
% (sin a, cos a) 이고, TikZ 의 수학각은 90 - a 다.
%   psi   = 30 deg  (수학각 60)   선수각
%   chi_r = 40 deg  (수학각 50)   대수 침로  -> beta_r = 10 deg
%   beta_c = 90 deg (수학각  0)   조류가 가는 방위 = 정동
% 길이: U_r = 42 mm, U_c = 14 mm
%   U_r 벡터 = 42 (sin40, cos40) = (27.00, 32.17)
%   U        = U_r + (14, 0)     = (41.00, 32.17),  |U| = 52.11 mm
%   chi = atan2(41.00, 32.17) = 51.88 deg  (수학각 38.12)
%   beta = chi - psi = 21.88 deg
\documentclass[border=5pt]{standalone}
% gnc-style.tex — 이 볼트의 TikZ 그림이 공유하는 스타일.
%
% 저널 그림 규약을 스타일 하나로 굳혀 둔다. 그림마다 선 굵기나 화살촉을
% 다시 정하지 않는다 — 그것이 그림이 제각각으로 보이는 원인이다.
%
% 쓰는 법:  % gnc-style.tex — 이 볼트의 TikZ 그림이 공유하는 스타일.
%
% 저널 그림 규약을 스타일 하나로 굳혀 둔다. 그림마다 선 굵기나 화살촉을
% 다시 정하지 않는다 — 그것이 그림이 제각각으로 보이는 원인이다.
%
% 쓰는 법:  % gnc-style.tex — 이 볼트의 TikZ 그림이 공유하는 스타일.
%
% 저널 그림 규약을 스타일 하나로 굳혀 둔다. 그림마다 선 굵기나 화살촉을
% 다시 정하지 않는다 — 그것이 그림이 제각각으로 보이는 원인이다.
%
% 쓰는 법:  \input{gnc-style}  (figures/src/ 안에서)

\usepackage{tikz}
\usepackage{amsmath}
\usepackage{xcolor}
\usetikzlibrary{arrows.meta, positioning, calc, fit, backgrounds, decorations.pathmorphing}

% ---- 색 --------------------------------------------------------------------
% 색은 강조 하나에만 쓴다. 나머지는 검정.
\definecolor{ink}{HTML}{1A1A1A}
\definecolor{accent}{HTML}{7C3AED}
\definecolor{warn}{HTML}{B91C1C}
\definecolor{soft}{HTML}{666666}

% 기하 그림(좌표계·유도)에서만 쓰는 넷. 블록선도는 위의 흑백 + accent 로 충분하다.
\definecolor{pathc}{HTML}{0D9488}   % 계획 경로 · 항적
\definecolor{errc}{HTML}{C2410C}    % 오차 (cross-track 등)
\definecolor{alng}{HTML}{B45309}    % along-track
\definecolor{flow}{HTML}{0369A1}    % 조류 · 유속

% ---- 선과 화살촉 -----------------------------------------------------------
\tikzset{
  % 신호선. 화살촉은 작고 뾰족한 것 하나뿐이다.
  sig/.style   = {ink, line width=0.5pt, -{Stealth[length=4pt, width=3pt]}},
  wire/.style  = {ink, line width=0.5pt},
  % 강조 경로 — 파선으로 구분한다. 흑백 인쇄에서도 살아야 하기 때문이다.
  asig/.style  = {accent, line width=0.5pt, densely dashed,
                  -{Stealth[length=4pt, width=3pt]}},
  awire/.style = {accent, line width=0.5pt, densely dashed},
  % 블록. 흰 바탕, 직각, 검은 테두리.
  blk/.style   = {draw=ink, line width=0.55pt, fill=white,
                  minimum width=17mm, minimum height=8mm, inner sep=2pt, font=\small},
  % 합산점
  sum/.style   = {draw=ink, line width=0.55pt, fill=white, circle,
                  inner sep=0pt, minimum size=4.6mm, font=\scriptsize},
  asum/.style  = {draw=accent, line width=0.55pt, fill=white, circle,
                  inner sep=0pt, minimum size=4.6mm, font=\scriptsize},
  % 분기점
  dot/.style   = {ink, fill=ink, circle, inner sep=0pt, minimum size=1.5mm},
  % 신호 이름 — 선 위에 작게
  sname/.style = {font=\small, inner sep=1.5pt},
  % 그림 안의 짧은 설명. 그림 안의 글은 최소로 둔다
  note/.style  = {font=\scriptsize, soft, inner sep=1.5pt},
  anote/.style = {font=\scriptsize, accent, inner sep=1.5pt},
  % 축
  axis/.style  = {ink, line width=0.5pt},
  tick/.style  = {font=\scriptsize, soft},
}

% 캡션 한 줄 — 그림 아래 가는 선과 함께
\newcommand{\gnccaption}[2]{%
  \node[anchor=north west, text width=#1, align=left, font=\scriptsize, ink]
        at (current bounding box.south west) {\vspace{2pt}#2};}
  (figures/src/ 안에서)

\usepackage{tikz}
\usepackage{amsmath}
\usepackage{xcolor}
\usetikzlibrary{arrows.meta, positioning, calc, fit, backgrounds, decorations.pathmorphing}

% ---- 색 --------------------------------------------------------------------
% 색은 강조 하나에만 쓴다. 나머지는 검정.
\definecolor{ink}{HTML}{1A1A1A}
\definecolor{accent}{HTML}{7C3AED}
\definecolor{warn}{HTML}{B91C1C}
\definecolor{soft}{HTML}{666666}

% 기하 그림(좌표계·유도)에서만 쓰는 넷. 블록선도는 위의 흑백 + accent 로 충분하다.
\definecolor{pathc}{HTML}{0D9488}   % 계획 경로 · 항적
\definecolor{errc}{HTML}{C2410C}    % 오차 (cross-track 등)
\definecolor{alng}{HTML}{B45309}    % along-track
\definecolor{flow}{HTML}{0369A1}    % 조류 · 유속

% ---- 선과 화살촉 -----------------------------------------------------------
\tikzset{
  % 신호선. 화살촉은 작고 뾰족한 것 하나뿐이다.
  sig/.style   = {ink, line width=0.5pt, -{Stealth[length=4pt, width=3pt]}},
  wire/.style  = {ink, line width=0.5pt},
  % 강조 경로 — 파선으로 구분한다. 흑백 인쇄에서도 살아야 하기 때문이다.
  asig/.style  = {accent, line width=0.5pt, densely dashed,
                  -{Stealth[length=4pt, width=3pt]}},
  awire/.style = {accent, line width=0.5pt, densely dashed},
  % 블록. 흰 바탕, 직각, 검은 테두리.
  blk/.style   = {draw=ink, line width=0.55pt, fill=white,
                  minimum width=17mm, minimum height=8mm, inner sep=2pt, font=\small},
  % 합산점
  sum/.style   = {draw=ink, line width=0.55pt, fill=white, circle,
                  inner sep=0pt, minimum size=4.6mm, font=\scriptsize},
  asum/.style  = {draw=accent, line width=0.55pt, fill=white, circle,
                  inner sep=0pt, minimum size=4.6mm, font=\scriptsize},
  % 분기점
  dot/.style   = {ink, fill=ink, circle, inner sep=0pt, minimum size=1.5mm},
  % 신호 이름 — 선 위에 작게
  sname/.style = {font=\small, inner sep=1.5pt},
  % 그림 안의 짧은 설명. 그림 안의 글은 최소로 둔다
  note/.style  = {font=\scriptsize, soft, inner sep=1.5pt},
  anote/.style = {font=\scriptsize, accent, inner sep=1.5pt},
  % 축
  axis/.style  = {ink, line width=0.5pt},
  tick/.style  = {font=\scriptsize, soft},
}

% 캡션 한 줄 — 그림 아래 가는 선과 함께
\newcommand{\gnccaption}[2]{%
  \node[anchor=north west, text width=#1, align=left, font=\scriptsize, ink]
        at (current bounding box.south west) {\vspace{2pt}#2};}
  (figures/src/ 안에서)

\usepackage{tikz}
\usepackage{amsmath}
\usepackage{xcolor}
\usetikzlibrary{arrows.meta, positioning, calc, fit, backgrounds, decorations.pathmorphing}

% ---- 색 --------------------------------------------------------------------
% 색은 강조 하나에만 쓴다. 나머지는 검정.
\definecolor{ink}{HTML}{1A1A1A}
\definecolor{accent}{HTML}{7C3AED}
\definecolor{warn}{HTML}{B91C1C}
\definecolor{soft}{HTML}{666666}

% 기하 그림(좌표계·유도)에서만 쓰는 넷. 블록선도는 위의 흑백 + accent 로 충분하다.
\definecolor{pathc}{HTML}{0D9488}   % 계획 경로 · 항적
\definecolor{errc}{HTML}{C2410C}    % 오차 (cross-track 등)
\definecolor{alng}{HTML}{B45309}    % along-track
\definecolor{flow}{HTML}{0369A1}    % 조류 · 유속

% ---- 선과 화살촉 -----------------------------------------------------------
\tikzset{
  % 신호선. 화살촉은 작고 뾰족한 것 하나뿐이다.
  sig/.style   = {ink, line width=0.5pt, -{Stealth[length=4pt, width=3pt]}},
  wire/.style  = {ink, line width=0.5pt},
  % 강조 경로 — 파선으로 구분한다. 흑백 인쇄에서도 살아야 하기 때문이다.
  asig/.style  = {accent, line width=0.5pt, densely dashed,
                  -{Stealth[length=4pt, width=3pt]}},
  awire/.style = {accent, line width=0.5pt, densely dashed},
  % 블록. 흰 바탕, 직각, 검은 테두리.
  blk/.style   = {draw=ink, line width=0.55pt, fill=white,
                  minimum width=17mm, minimum height=8mm, inner sep=2pt, font=\small},
  % 합산점
  sum/.style   = {draw=ink, line width=0.55pt, fill=white, circle,
                  inner sep=0pt, minimum size=4.6mm, font=\scriptsize},
  asum/.style  = {draw=accent, line width=0.55pt, fill=white, circle,
                  inner sep=0pt, minimum size=4.6mm, font=\scriptsize},
  % 분기점
  dot/.style   = {ink, fill=ink, circle, inner sep=0pt, minimum size=1.5mm},
  % 신호 이름 — 선 위에 작게
  sname/.style = {font=\small, inner sep=1.5pt},
  % 그림 안의 짧은 설명. 그림 안의 글은 최소로 둔다
  note/.style  = {font=\scriptsize, soft, inner sep=1.5pt},
  anote/.style = {font=\scriptsize, accent, inner sep=1.5pt},
  % 축
  axis/.style  = {ink, line width=0.5pt},
  tick/.style  = {font=\scriptsize, soft},
}

% 캡션 한 줄 — 그림 아래 가는 선과 함께
\newcommand{\gnccaption}[2]{%
  \node[anchor=north west, text width=#1, align=left, font=\scriptsize, ink]
        at (current bounding box.south west) {\vspace{2pt}#2};}


\tikzset{
  ttl/.style   = {font=\small\bfseries, ink},
  blurb/.style = {font=\scriptsize, ink, align=left},
  ang/.style   = {soft, line width=0.45pt},
  vec/.style   = {line width=0.9pt, -{Stealth[length=4.6pt,width=3.4pt]}},
  thin/.style  = {soft, line width=0.35pt, densely dashed},
}

\definecolor{relc}{HTML}{0D9488}   % 대수 — pathc 와 같은 청록
\definecolor{gndc}{HTML}{1D4ED8}   % 대지 — 파랑

\begin{document}
\begin{tikzpicture}[x=1mm, y=1mm]

% =====================================================================
% 위 — 해류 삼각형
% =====================================================================
\begin{scope}[shift={(0,44)}]
  \node[ttl, anchor=north west] at (-10,62) {the current triangle};

  % ---- 축 ----
  \draw[thin] (0,-4) -- (0,52);
  \node[font=\small, anchor=south west] at (1.0,51.5) {$x$~(North)};
  \draw[axis, -{Stealth[length=4pt,width=3pt]}] (-4,0) -- (58,0);
  \node[font=\small, anchor=west] at (58.8,0) {$y$~(East)};

  % ---- 선체 : 선수는 방위각 30 도 ----
  \draw[ink, line width=0.8pt, -{Stealth[length=4pt,width=3pt]}] (0,0) -- (24.0,41.6);
  \node[ink, font=\small, anchor=south east] at (23.0,41.8) {$x_b$};
  \draw[soft, line width=0.6pt, -{Stealth[length=4pt,width=3pt]}] (0,0) -- (17.3,-10.0);
  \node[soft, font=\small, anchor=north west] at (17.0,-10.4) {$y_b$};
  %  선체는 작게. 크면 원점의 각들을 덮는다
  \fill[ink] (1.9,3.2) -- (-2.9,1.6) -- (-1.2,-3.2) -- cycle;

  % ---- 세 속도 벡터 ----
  \draw[relc, vec] (0,0) -- (27.00,32.17);
  \node[relc, font=\small, anchor=south east] at (23.4,33.0) {$\mathbf{U}_r$};
  \draw[flow, vec] (27.00,32.17) -- (41.00,32.17);
  \node[flow, font=\small, anchor=south] at (34.0,33.0) {$\mathbf{U}_c$};
  \draw[gndc, vec] (0,0) -- (41.00,32.17);
  \node[gndc, font=\small, anchor=west] at (41.8,31.6) {$\mathbf{U}$};

  % ---- 각 ----
  %  세 침로각은 모두 북에서 시작하므로, 이름표를 **arc 가 시작하는 자리**에
  %  나란히 둔다. 중앙에 두면 벡터 셋과 겹친다
  \draw[ang] (0,22) arc (90:60:22);
  \node[soft, font=\small, anchor=south west] at (1.4,22.0) {$\psi$};
  \draw[relc, line width=0.45pt] (0,28) arc (90:50:28);
  \node[relc, font=\small, anchor=south west] at (1.4,28.0) {$\chi_r$};
  \draw[gndc, line width=0.45pt] (0,34) arc (90:38.12:34);
  \node[gndc, font=\small, anchor=south west] at (1.4,34.0) {$\chi$};

  \draw[accent, line width=0.45pt] (7.00,12.12) arc (60:38.12:14);
  \node[accent, font=\small] at (14.8,7.6) {$\beta$};
  \draw[relc, line width=0.45pt] (9.00,15.59) arc (60:50:18);
  \node[relc, font=\small] at (14.6,15.6) {$\beta_r$};

  % ---- 조류가 가는 방위 : 조류 화살표가 시작하는 곳에서 북에서 잰다 ----
  \draw[thin] (27.00,32.17) -- (27.00,46.0);
  \draw[flow, line width=0.45pt] (27.00,41.17) arc (90:0:9);
  \node[flow, font=\small] at (31.6,40.4) {$\beta_c$};

  \node[note, anchor=north west, text width=60mm] at (-10,-14) {%
    Drawn with $\psi = 30^\circ$, $\chi_r = 40^\circ$ and a current running due
    east, which puts $\chi$ at $51.9^\circ$ and the crab angle at
    $\beta = 21.9^\circ$.};
\end{scope}

% =====================================================================
% 아래 — 같은 대지속도를 선체 축에서 보면 u 와 v
% =====================================================================
\begin{scope}[shift={(6,-26)}]
  %  제목은 x_b 이름표(0,30.6)보다 위에 둔다 — 같은 높이면 글자가 겹친다
  \node[ttl, anchor=north west] at (-16,40) {the same $\mathbf{U}$, read on the hull's own axes};

  \draw[axis, -{Stealth[length=4pt,width=3pt]}] (0,0) -- (0,30);
  \node[ink, font=\small, anchor=south] at (0,30.6) {$x_b$};
  \draw[axis, -{Stealth[length=4pt,width=3pt]}] (0,0) -- (32,0);
  \node[ink, font=\small, anchor=west] at (32.8,0) {$y_b$};

  \draw[gndc, vec] (0,0) -- (11.0,25.0);
  \node[gndc, font=\small, anchor=south west] at (11.2,25.0) {$\mathbf{U}$};
  \draw[alng, line width=0.8pt, -{Stealth[length=4pt,width=3pt]}] (0,0) -- (0,25.0);
  \node[alng, font=\small, anchor=east] at (-1.2,13.0) {$u$};
  \draw[errc, line width=0.8pt, -{Stealth[length=4pt,width=3pt]}] (0,25.0) -- (11.0,25.0);
  \node[errc, font=\small, anchor=north] at (5.5,24.0) {$v$};
  \draw[soft, line width=0.35pt, densely dotted] (11.0,25.0) -- (11.0,0);
  \draw[accent, line width=0.45pt] (0,9) arc (90:66.2:9);
  \node[accent, font=\small] at (4.6,8.0) {$\beta$};

  \node[note, anchor=north west, text width=60mm] at (-16,-6) {%
    $u$ is surge and $v$ is sway, both taken \emph{over the ground}. The
    water-relative pair $u_r$, $v_r$ is read on these same two axes, and
    $\beta_r$ is the angle it makes with $x_b$.};
\end{scope}

% =====================================================================
% 오른쪽 — 다섯 각의 이름, 그리고 그 사이의 식
% =====================================================================
\begin{scope}[shift={(86,0)}]
  \node[ttl, anchor=north west] at (0,102) {five angles, and what separates them};

  \node[blurb, anchor=north west, text width=84mm] at (0,95) {%
    \textcolor{soft}{$\psi$}\quad \textbf{heading} --- the direction the bow points\\[2pt]
    \textcolor{gndc}{$\chi$}\quad \textbf{course over ground} --- the direction the
      vessel travels \emph{over the ground}\\[2pt]
    \textcolor{relc}{$\chi_r$}\quad \textbf{course through the water} --- the direction
      it travels \emph{through the water}\\[2pt]
    \textcolor{accent}{$\beta$}\quad \textbf{crab angle} --- the gap between the bow and
      the course over ground\\[2pt]
    \textcolor{relc}{$\beta_r$}\quad \textbf{drift angle} --- the gap between the bow and
      the course through the water\\[2pt]
    \textcolor{flow}{$\beta_c$}\quad the direction the current flows \emph{towards},
      measured from north};

  \draw[soft, line width=0.3pt] (0,38) -- (84,38);

  \node[anchor=north west, font=\small] at (0,35)
       {$\mathbf{U} = \mathbf{U}_r + \mathbf{U}_c$};
  \node[blurb, anchor=north west, text width=50mm] at (34,35) {%
    speed over ground $=$ speed through the water $+$ current};
  \node[anchor=north west, font=\small] at (0,26)
       {$U = \sqrt{u^2+v^2}, \qquad U_r = \sqrt{u_r^2+v_r^2}$};
  \node[anchor=north west, font=\small] at (0,16)
       {$\chi = \psi + \beta, \qquad \beta = \operatorname{atan2}(v,\ u)$};
  \node[anchor=north west, font=\small] at (0,6)
       {$\chi_r = \psi + \beta_r, \qquad \beta_r = \arcsin(v_r/U_r)$};

  \draw[soft, line width=0.3pt] (0,-2) -- (84,-2);

  \node[blurb, anchor=north west, text width=84mm, soft] at (0,-5) {%
    \textbf{What measures what.} A compass gives $\psi$. GNSS gives $\chi$ and $U$,
    by differencing position. Nothing gives $\beta$: it is what is left over,
    $\chi-\psi$.};
  \node[blurb, anchor=north west, text width=84mm, gndc] at (0,-21) {%
    With no current $\mathbf{U} = \mathbf{U}_r$, so $\chi = \chi_r$ and
    $\beta = \beta_r$. Everything that separates the two triangles in this figure
    is the current, and nothing else.};
  \node[blurb, anchor=north west, text width=84mm, warn] at (0,-37) {%
    An autopilot holds $\psi$. The path is walked along $\chi$. Hold the first and
    the vessel travels sideways off the path by $\beta$ --- which is the error
    §5-5 has to remove.};
\end{scope}

\end{tikzpicture}
\end{document}
